\subsection{Contraction of arrows of a groupoid}

Let G be a groupoid and let F be a subset of the set of arrows of G. Let H be the distinguished subgroupoid of G generated by F. The groupoid G/H is called the groupoid obtained from G by contracting the set of arrows F and is denoted G/F. If p denotes the canonical morphism from G to G/F, we have \(p(o(f)) = p(t(f))\) and \(p(f) = e_{p(o(f))}\) for every arrow \(f \in F\) .

This morphism satisfies the following universal property.

PROPOSITION 7. --- For any morphism of groupoids \(\varphi\colon G\to G'\) which maps every arrow belonging to F to an identity element of \(G'\) , there exists a unique morphism of groupoids \(\varphi'\colon G/F\to G'\) such that \(\varphi'\circ p=\varphi\) .

The kernel of \(\varphi\) is a distinguished subgroupoid of G, and the set of its arrows contains F. By definition, H is the smallest distinguished subgroupoid of G whose set of arrows contains F; it is therefore contained in the kernel of \(\varphi\) . The proposition then follows from II, p. 170, prop. 3.

Denote by \(\Gamma\) the subquiver of G with vertex set Som(G) and arrow set F. The orbits of H are identified with the connected components of the quiver \(\Gamma\) . By definition of the quotient of a groupoid by a distinguished subgroupoid, the vertex set of G/F is identified with the set of connected components of \(\Gamma\) ; in other words, it is the quotient set of Som(G) by the finest equivalence relation such that \(o(f)\) is equivalent to \(t(f)\) for every arrow \(f \in \mathrm{F}\) .

The map p induces, by passing to quotients, a bijection from the set of orbits of G to the set of orbits of G/F (II, p. 170, remark 1).

The purpose of this section is to compute the homomorphisms induced by p by passing to the isotropy groups, which, according to Prop. 2 (II, p. 170), amounts to computing the isotropy groups of the subgroupoid H.

PROPOSITION 8. --- Let δ be the canonical morphism from the free groupoid \(\mathrm{Grp}(\Gamma)\) into G extending the canonical injective morphism from Γ into G. Let a be a vertex of G and let A be the orbit of a in G. For every b ∈ A, let \(f_{ab}\) be an arrow of G connecting a to b. The isotropy group \(H_{a}\) is the distinguished subgroup of \(G_{a}\) generated by the elements \(\operatorname{Int}(f_{ab})(\delta(c))\) , where b ranges over the set A and c ranges over the elements of \(\mathrm{Grp}(\Gamma)_{b}\) .

If x and y are vertices of G belonging to the same orbit, with \(x \neq a\) ; we further fix an arrow \(f_{xy}\) of G that connects x to y. For every point x of G, let \(N_{x}\) be the distinguished subgroup of \(G_{x}\) generated by the elements \(\operatorname{Int}(f_{xy})(\delta(c))\) , where y ranges over the orbit of x in G and c ranges over the group \(\operatorname{Grp}(\Gamma)_{y}\) . For all \(f \in G_{xy}\) and every \(g \in N_{y}\) , we have \(fgf^{-1} \in N_{x}\) . Let \(x \in \operatorname{Som}(G)\) . We have \(\delta(\operatorname{Grp}(\Gamma)_{x}) \subset N_{x}\) , by definition of \(N_{x}\) . By definition of the groupoid H, \(\delta(f)\) is an arrow of H for every arrow f of \(\operatorname{Grp}(\Gamma)\) ; consequently, \(\delta(\operatorname{Grp}(\Gamma)_{x})\) is contained in \(H_{x}\) . Since H is a distinguished subgroupoid of G, we then have

\[
\operatorname{Int} (f _ {x y}) (\delta (c)) \in \operatorname{Int} (f _ {x y}) (\mathrm{N} _ {y}) \subset \operatorname{Int} (f _ {x y}) (\mathrm{H} _ {y}) \subset \mathrm{H} _ {x}
\]

for all \(y \in \operatorname{Som}(\mathbf{G})\) , then \(\mathrm{N}_x \subset \mathrm{H}_x\) .

Let \(x\) and \(y\) be vertices of \(\Gamma\) . Set \(\mathrm{N}_{x,y} = \varnothing\) if \(x\) and \(y\) do not belong to the same connected component of \(\Gamma\) . Otherwise, let \(c\) and \(c'\) be path classes connecting \(x\) to \(y\) in the graph \(\widetilde{\Gamma}\) associated with \(\Gamma\) . Let \(z\) be a vertex of G belonging to the orbit of \(y\) and let \(\ell\) be an element of \(\mathrm{Grp}(\Gamma)_z\) ; in the group \(\mathrm{G}_x\) one has the equality:

\[
\begin{array}{r l} & {\delta (c) f _ {y z} \delta (\ell) f _ {y z} ^ {- 1} \delta (c ^ {\prime}) ^ {- 1}} \\ & {\qquad = \delta (c) f _ {y z} f _ {x z} ^ {- 1} \big (f _ {x z} \delta (\ell) f _ {x z} ^ {- 1} \big) f _ {x z} f _ {y z} ^ {- 1} \delta (c) ^ {- 1} \delta (c (c ^ {\prime}) ^ {- 1}).} \end{array}
\]

The arrow \(\delta(c(c')^{-1})\) is the class of a loop at \(x\) in the quiver \(\Gamma\) , hence belongs to \(\mathrm{N}_x\) , just as the arrow \(f_{xz}\delta(\ell)f_{xz}^{-1}\) . Since \(\delta(c)f_{yz}f_{xz}^{-1}\) belongs to \(\mathrm{G}_x\) and \(\mathrm{N}_x\) is a distinguished subgroup of \(\mathrm{G}_x\) , it follows that \(\delta(c)f_{yz}\delta(\ell)f_{yz}^{-1}\delta(c')^{-1}\) belongs to \(\mathrm{N}_x\) . Consequently, \(\delta(c)\mathrm{N}_y\subset \mathrm{N}_x\delta(c')\) . By symmetry, we have \(\delta(c)\mathrm{N}_y = \mathrm{N}_x\delta(c')\) . This implies that this set does not depend on the choices of \(c\) and \(c'\) ; it is denoted \(\mathrm{N}_{x,y}\) .

Let N be the subquiver of G whose set of vertices is Som(G) and whose set of arrows from x to y is equal to \(N_{x,y}\) for every pair \((x,y)\) of vertices of G. It is a distinguished subgroupoid of G whose set of arrows contains F; consequently, H is a subgroupoid of N. In particular, \(H_{a} \subset N_{a}\) , whence the equality.

COROLLARY 1. --- Suppose further that, for every arrow f of F, we have \(o(f) = t(f)\) . Then, the origin and the target of every path in the quiver \(\Gamma\) coincide. Let a be a vertex of G and let A be its orbit in G. For every element c of F whose origin b belongs to A, fix an arrow \(f_{c}\) of G connecting a to b, and set \(\kappa(c) = \operatorname{Int}(f_{c})(\delta(c)) = f_{c}\delta(c)f_{c}^{-1}\) .

The group \(H_{a}\) is then the distinguished subgroup of \(G_{a}\) generated by the arrows \(\kappa(c)\) .

By prop. 8, the arrows \(\kappa(c)\) are contained in the group \(H_a\) . Let \(N_a\) be the smallest normal subgroup of \(G_a\) containing them. By this proposition, it suffices to prove that, for every vertex \(b\) of G belonging to A, every arrow \(f_{ab}\) connecting \(a\) to \(b\) in G and every element \(c\) of \(\mathrm{Grp}(\Gamma)_b\) , \(\operatorname{Int}(f_{ab})(\delta(c))\) belongs to \(N_a\) . Since the origin and the target of each element of F are equal, a loop at a vertex \(b\) in the graph \(\widetilde{\Gamma}\) is of the form \((b,(f_1,\varepsilon_1),b,\ldots,(f_n,\varepsilon_n),b)\), where for every \(i\in\{1,\ldots,n\}\), \(f_i\) is an element of F with origin \(b\) and \(\varepsilon_i\in\{-1,1\}\). We have

\[
\operatorname{Int} (f _ {a b}) (\delta (c)) = f _ {a b} f _ {c} ^ {- 1} \kappa (f _ {1}) ^ {\varepsilon_ {1}} \dots \kappa (f _ {n}) ^ {\varepsilon_ {n}} f _ {c} f _ {a b} ^ {- 1}.
\]

Since \(f_{ab}f_c^{-1} \in \mathrm{G}_a\) , it follows that \(\operatorname{Int}(f_{ab})(\delta(c))\) belongs to \(\mathrm{N}_a\) .

COROLLARY 2. --- If the graph associated with the quiver \(\Gamma\) is a forest, the homomorphism \(p_a\) from \(\mathrm{G}_a\) to \((\mathrm{G} / \mathrm{F})_{p(a)}\) is bijective, for every vertex a of G.

Indeed, under this hypothesis, we have \(\mathrm{Grp}(\Gamma)_{b}=\{e_{b}\}\) for every vertex b of G (corollary 1 of II, p. 173). It then follows from Proposition 8 that, for every vertex a of G, the group \(H_{a}\) is reduced to the identity element.

